\documentclass[12pt,a4paper]{article}

\usepackage[margin=1in]{geometry}
\usepackage{amsmath}
\usepackage[table]{xcolor}
\usepackage{graphicx}
\usepackage{hyperref}

\title{Technical Writing: LaTeX Demonstration}
\author{Aneek Chakraborty}
\date{\today}

\begin{document}

\maketitle
\tableofcontents
\newpage

\section{Introduction}

LaTeX is a document preparation system used for
structured technical and scientific writing.

This document demonstrates \textbf{bold text},
\textit{italic text}, and
\textcolor{blue}{coloured text}.

\subsection{Learning Objectives}

\begin{itemize}
    \item Organise a document using sections.
    \item Typeset mathematical expressions.
    \item Create tables and figures.
    \item Use automatic cross-references.
\end{itemize}

\section{Mathematical Expressions}

Inline mathematics fits within a sentence:
\( a^2+b^2=c^2 \).

A numbered equation appears separately:

\begin{equation}
    \sum_{i=1}^{n} i = \frac{n(n+1)}{2}
    \label{eq:sum}
\end{equation}

Equation~\eqref{eq:sum} gives the sum of the first
\(n\) positive integers.

A matrix:

\[
A =
\begin{bmatrix}
1 & 2 \\
3 & 4
\end{bmatrix}
\]

A piecewise function:

\[
f(x)=
\begin{cases}
-x, & x<0,\\
x, & x\geq 0.
\end{cases}
\]

\section{Tables}

Table~\ref{tab:tools} compares three tools.

\begin{table}[htbp]
    \centering
    \caption{Technical Writing Tools}
    \label{tab:tools}

    \begin{tabular}{|l|p{7cm}|}
        \hline
        \rowcolor{blue!15}
        \textbf{Tool} & \textbf{Purpose} \\
        \hline
        LaTeX & Preparing technical documents \\
        \hline
        Beamer & Creating presentations \\
        \hline
        Doxygen & Generating source-code documentation \\
        \hline
    \end{tabular}
\end{table}

\section{Figures}

Figure~\ref{fig:workflow} is drawn using the
picture environment, without an external image.

\begin{figure}[htbp]
    \centering
    \setlength{\unitlength}{1mm}
    \begin{picture}(110,25)
        \put(0,5){\framebox(30,15){Source}}
        \put(30,12){\vector(1,0){10}}
        \put(40,5){\framebox(30,15){Compile}}
        \put(70,12){\vector(1,0){10}}
        \put(80,5){\framebox(30,15){PDF}}
    \end{picture}
    \caption{LaTeX compilation workflow}
    \label{fig:workflow}
\end{figure}

\section{Compilation Procedure}

\begin{enumerate}
    \item Write and save the LaTeX source.
    \item Compile using pdfLaTeX.
    \item Check the output for errors.
    \item Compile again if references need updating.
\end{enumerate}

\noindent
\fbox{\parbox{0.9\linewidth}{
    \textbf{Key point:}
    The source file contains commands and content.
    The compiled PDF displays the formatted document.
}}

\section{Conclusion}

LaTeX provides consistent formatting and automatic
numbering for structured documents.
The LaTeX Project provides further information
\cite{latexproject}.

\begin{thebibliography}{9}
    \bibitem{latexproject}
    The LaTeX Project.
    \url{https://www.latex-project.org/}
\end{thebibliography}

\end{document}
